% Part III: waveform recordings. Sources: docs/wfdb-collection.md (every section);
% docs/adapters.md (what add-ons do; the protocol is chap:addons); docs/importing.md "Files";
% docs/administration.md "Files and recordings", "Linking recordings", "Identifiers in
% recordings"; docs/edf.md (users' part); viewers/vitalstream/docs/data-model.md. Checked against
% libs/engine/src/sources/wfdb_folder.rs and sources/mod.rs, libs/engine/src/index.rs
% (cache_key), medorama/src/tables.rs (files of rows), recordings.rs (find, add, check),
% linkage.rs (relink), jobs.rs (automatic indexing), web/js/views/recordings.js, encounters.js.
\chapter{Waveform recordings}
\label{chap:recordings}

This chapter describes the waveform recordings that Medorama reads, how they come into the
catalog, how they find their encounters, and how they are indexed. It is written for
administrators and data managers.

A \emph{recording} is a long waveform recording of a patient: the bedside monitors' signals of an
ICU stay, an ambulatory ECG (a Holter), an EEG, or a sleep study (polysomnography). Medorama
reads WFDB records, folders of records and WFDB collections, EDF, EDF+, BDF and BDF+ files, and
the formats that installed add-ons read. Recordings are shown in VitalStream
(\cref{chap:vitalstream}), from an \emph{index} that the catalog builds once, and are linked to
their patients' encounters. The patient's page shows them in the ICU, Holter, EEG or Sleep column
of the Patient Journey. A recording, like every file, comes in only as a row of a table that names
its patient (\cref{chap:importing}).

\section{What a recording becomes}
\label{sec:recordings-model}

Whatever its format, a recording is read into the same parts (\cref{tab:recordings-model}).
Three principles hold throughout:
\begin{itemize}
\item \textbf{Only what the recording states.} The readers do no resampling, interpolation, gap
  filling, re-timing or deduplication: every sample drawn is a stored sample.
\item \textbf{Ambiguity is shown, not resolved.} Where the recording does not settle something,
  such as a break's exact length or a start time, the viewer says so.
\item \textbf{Times are UTC}, as Unix seconds. The time zone shown is a setting of the viewer
  only.
\end{itemize}

\begin{table}[htbp]
\caption{What a recording is made of.}
\label{tab:recordings-model}
\begin{tabularx}{\linewidth}{@{}P{2.4cm}L@{}}
\toprule
{\tabhead Term} & {\tabhead Meaning} \\
\midrule
Channel & One signal as the user sees it (ECG lead II, ART, Pleth \dots), defined by a
  modality, a label and a unit. \\
Piece & Samples of one channel that the recording proves contiguous, at one constant sampling
  rate: sample $i$ is at $t_0 + i/f_s$. Pieces are drawn apart, never joined. \\
Lane & The pieces of a channel never overlap in time. A signal recorded twice at the same time
  goes into a second lane, which is a channel of its own (\cref{sec:recordings-overlap}). \\
Time not stated & A piece for which the recording states no time. It is placed next to its
  neighbors and marked as such. \\
Break & A place where contiguity is not proven, with its reason, its length when known, whether
  that length is exact, how it was measured, and the channels it concerns. \\
Vital-sign series & Values at their times, by name and unit. There is one series per name and
  source; series are never merged. \\
Alarm & An interval with a text, a kind and a level (1 advisory, 2 warning, 3 crisis). Alarm
  colors are the only red on the viewers' pages. EDF+ annotations are events of the kind
  \code{annotation}. \\
Marker & A mark at a sample of a piece (an R wave, a pacing spike), by name. \\
\bottomrule
\end{tabularx}
\end{table}

Nothing is assumed that a recording does not state. A WFDB record whose header gives no sampling
frequency is not read. A signal with a gain of 0 (uncalibrated) is shown in stored counts, with
the unit \code{adu}. A signal whose header gives no units has none: WFDB's old defaults of
250\,Hz, a gain of 200 and millivolts are not used. Missing samples are drawn as a gap in the
line, never as zero, and zero is a value, never ``missing''.

\section{Bringing recordings in}
\label{sec:recordings-in}

A recording comes in as the \code{file} of a table's row. The file becomes a recording when the
row's or the table's category is a \textbf{recording category} (\code{icu}, \code{holter},
\code{eeg} or \code{sleep}), or when neither the row nor the table names a category
(\cref{chap:importing}). A recording without a category goes to the ICU column. The file may be
one of the following:
\begin{itemize}
\item \textbf{A WFDB header} (\file{.hea}). The recording is that record only, or a
  multi-segment record whole (\cref{sec:recordings-multi}).
\item \textbf{A folder.} A \emph{WFDB folder}, one that holds WFDB records itself (a
  \file{.hea} file or a \file{RECORDS} file) or an \file{index.csv}, is itself the recording,
  taken whole with its subfolders, within the limits of \cref{sec:recordings-limits}. Any other
  folder adds only the recordings directly inside it (WFDB records, WFDB folders, EDF files and
  files an add-on reads); a subfolder that is not itself a WFDB folder is not searched.
\item \textbf{An EDF, EDF+, BDF or BDF+ file} (\file{.edf}, \file{.bdf};
  \cref{sec:recordings-edf}).
\item \textbf{A file of a format that an installed add-on reads}, such as BedMasterEx
  \file{.Stp} (\cref{sec:recordings-addons}).
\end{itemize}

The following rules keep recordings safe:
\begin{itemize}
\item \textbf{Recordings are only read.} The catalog stores paths, never copies. It never writes,
  moves or deletes a recording, and never writes into a recording's folder (the folder that holds
  a recording file, or anywhere inside a WFDB folder). A recording that lies where the catalog
  writes (its results, indexes or home folder), or that holds one of them, is refused at import.
\item \textbf{Inside \code{data\_roots}.} When the setting names folders, a recording outside
  them is left out with a warning.
\item \textbf{Once.} Paths are stored with symbolic links resolved, so a recording is never in
  the catalog twice. A second table that names it claims it too, and the recording stays while
  either table does.
\item \textbf{Hidden entries} (names that start with \code{.}) are skipped when a folder is
  searched.
\item \textbf{WFDB recordings need WFDB support.} A Medorama built without it
  (\cref{sec:installation-no-wfdb}) skips each WFDB record, folder of records or collection it
  finds, and lists it with the reason ``this build has no WFDB support''.
\end{itemize}
A recording takes its row's category (ICU by default), and is linked to the row's encounter (\cref{sec:recordings-linking}), indexed in the background
(\cref{sec:recordings-indexing}), and listed in \ui{Import} $\rightarrow$ \ui{Recordings}
(\cref{sec:recordings-tab}).

\screenshot{patient-holter.png}{A patient's page with a Holter recording in its column of the Patient Journey.}{fig:recordings-holter}

\section{WFDB records}
\label{sec:recordings-wfdb}

Medorama and the three VitalStream viewers read WFDB records themselves, with any signal format
the WFDB library reads (through wfdb-x). They read the same recording to the same index, byte for
byte. A recording may be one record, a multi-segment record, a folder of records read as one
session, or a WFDB collection (\cref{tab:recordings-forms}).

\begin{table}[htbp]
\caption{The forms of WFDB recordings.}
\label{tab:recordings-forms}
\begin{tabularx}{\linewidth}{@{}P{3.4cm}LP{3.2cm}@{}}
\toprule
{\tabhead Form} & {\tabhead What it is} & {\tabhead Format shown} \\
\midrule
One record & A header file (\file{rec.hea}) and the signal files it names &
  \code{WFDB record} \\
A multi-segment record & A master header (\code{name/n} on its record line) and its segments &
  \code{WFDB multi-segment record} \\
A folder of records & A folder without \file{index.csv}: every record in it and in its
  subfolders, on one time axis & \code{WFDB records} \\
A collection & A folder with an \file{index.csv} that lists its records, and CSV files that say
  how they fit together & \code{WFDB collection} \\
\bottomrule
\end{tabularx}
\end{table}

\subsection{One record}
\label{sec:recordings-one}

A record's header file (\file{rec.hea}), opened in a viewer or named in a table, opens
\emph{that record only}: its signals, read from the signal files its header names, with its
\file{.mrk} markers. Its start is the header's base date and time, taken as UTC; a header that
states no date and time gives data marked \emph{time not stated}, which starts at 0
(1970-01-01). A record whose header states no sampling frequency, or one that is not a finite
number greater than 0 (\code{nan}), is not read.

\paragraph{Samples stated and samples held.} A header may state more samples than its signal
files hold. The number of samples read is the least of the number the header states and, for
each signal file whose signals all have a storage format of fixed size, the frames that file can
hold: $\lceil 8 \times \text{file size in bytes} / \text{bits per frame} \rceil$, where a frame
takes, for each signal of that file, its samples per frame (at least 1) times the bits of its
format (8 and 80: 8; 16, 61 and 160: 16; 212: 12; 310 and 311: 32/3; 24: 24; 32: 32). Files of
other formats (compressed) and files that do not exist are not counted. Reading also stops at
the end of the data, however many samples are stated. The file's byte offset is not subtracted,
so a well-formed record keeps the number its header states.

\paragraph{Files the record reads} are the header, the signal files it names (an absolute name as
given, a name starting with \code{\textasciitilde} in the home folder, any other next to the
header), its \file{.mrk} file and, for a multi-segment record, those of its segments. A viewer
opens none of them when it leads outside its data roots, symbolic links resolved
(\cref{sec:formats-reference-roots}).

\subsection{Multi-segment records}
\label{sec:recordings-multi}

A multi-segment record (WFDB's \code{name/n} record line, followed by $n$ segment lines) is read
as one record.
\begin{itemize}
\item Its \textbf{segments} follow each other from the record's start (the master header's base
  date and time), each at its offset: the samples of the segments before it, divided by the
  record's sampling frequency. Each signal of each segment is a piece of the channel of its name.
\item A \textbf{null segment} (\code{\textasciitilde}) is a break of known length (its samples
  divided by the sampling frequency), shown with the reason \code{null segment}.
\item A \textbf{layout segment} (the first segment, of length 0, of a record with a variable
  layout) names the record's signals and holds no samples. In a variable layout, a segment holds
  some of the signals; a signal absent from a segment is a break in its channel only.
\item The record's \textbf{markers} are the master's \file{.mrk} file: sample numbers from the
  record's start, and \code{chan} numbers of the master header's (the layout's) signals. Each
  marker goes to the segment its sample falls in, and to the signal of that name there; a marker
  in a null segment is left out.
\item Opening one \textbf{segment's own header} (\file{rec\_0001.hea}) reads that segment
  alone, as one record.
\end{itemize}
Segments are found next to their master header. Medorama and the viewers index a multi-segment
record as one recording, whose files are the master header, the layout header, and every
segment's header, signal files and markers (they make its index key,
\cref{sec:formats-reference-key}).

\subsection{A folder of records: one session}
\label{sec:recordings-session}

A folder without \file{index.csv} is read as \emph{one session}: every record in it and in its
subfolders, on one time axis. This is how a patient's records, split into folders by device or by
day, are seen together.
\begin{enumerate}
\item \textbf{Records} are the \file{.hea} files below the folder, at every level, in the order
  of their path relative to the folder (\code{/} between folders, compared character by
  character). Hidden entries and symbolic links are skipped. The segments and layout header of a
  multi-segment record are part of that record, never records of their own.
\item \textbf{Placement.} A record whose header states a base date and time starts there (UTC).
  A record without one is placed after the data before it, in that order, marked \emph{time not
  stated}.
\item A record that cannot be read, whose header states no valid sampling frequency, or one of
  whose files leads outside the data roots, is left out; the recording's notes say how many were
  left out.
\item A \file{RECORDS} file is not needed and not read. \file{vitals.csv}, \file{events.csv} and
  \file{gaps.csv}, when present in the folder itself, are read as for a collection.
\end{enumerate}
When a tree of folders is searched for recordings (Medorama's import), a folder is a recording
when it holds records itself (a \file{.hea} file or a \file{RECORDS} file) or an
\file{index.csv}; it is then taken whole, its subfolders included. The viewers' \ui{Open multiple
records\dots} opens any folder.

\screenshot{vitalstream-session.png}{A folder of records read as one session in VitalStream: a record without a start time, marked time not stated, and a signal recorded twice in a second lane, II (2).}{fig:recordings-session}

\subsection{WFDB collections}
\label{sec:recordings-collection}

A \emph{WFDB collection} is a folder of records with a few CSV files beside them, which say how
the records fit together in time and what else the recording held: vital signs, alarms, breaks.
Converters write collections (the STP add-on's \code{wfdb} export is one), and anyone can write
one with a few lines of code.
\begin{lstlisting}[style=text]
recording/                      the collection (one recording)
|-- index.csv                   the records, in order (required: it makes the folder a collection)
|-- gaps.csv                    breaks between records (optional)
|-- vitals.csv                  vital-sign values (optional)
|-- events.csv                  alarms and other events with a duration (optional)
|-- README.txt                  anything for people (ignored by the readers)
|-- ecg/                        a folder per modality, each with its RECORDS file
|   |-- RECORDS
|   |-- ecg_250_0001.hea  ecg_250_0001.dat  ecg_250_0001.mrk
|   `-- ...
`-- pressure/
    `-- ...
\end{lstlisting}
\begin{itemize}
\item Records are standard WFDB records; each signal of each record becomes a piece of a channel.
\item Every record's header states its sampling frequency: the readers assume none, so a
  collection with a record that does not is not read.
\item Records may sit in subfolders; \file{index.csv} names them by their path relative to the
  collection, with \code{/} between folders (\code{ecg/ecg\_250\_0001}).
\item The CSV files are UTF-8, comma-separated, with a header row, as Python's \code{csv} module
  writes them (quotes around fields that need them). Columns not described here are ignored, so a
  writer may add its own.
\item Times are Unix seconds (UTC), written as decimal numbers.
\item A collection lists its records itself, and is not subject to the limits of a folder
  (\cref{sec:recordings-limits}).
\end{itemize}

\begin{table}[htbp]
\caption{\file{index.csv}: one row per record, in the order the records were recorded.}
\label{tab:recordings-index}
\begin{tabularx}{\linewidth}{@{}P{2.4cm}P{1.6cm}L@{}}
\toprule
{\tabhead Column} & {\tabhead Required} & {\tabhead Meaning} \\
\midrule
\code{record} & Yes & The record's path relative to the collection, without extension. It
  stays inside the collection: a name that is absolute (starts with \code{/}, \code{\textbackslash}
  or \code{\textasciitilde}, or has a drive) or has a \code{..} component is not read \\
\code{modality} & & The kind of signals (\code{ecg}, \code{pressure}, \code{resp},
  \code{spo2} \dots); channels are grouped by it. Default \code{wfdb} \\
\code{start\_unix} & & The time of the record's first sample; empty: the record states no time
  (placed as below) \\
\code{stretch} & & A label shared by the records of one stretch of contiguous data; records
  without a start time that share it are placed together \\
\bottomrule
\end{tabularx}
\end{table}
Other columns of \file{index.csv} (\code{fs}, \code{n\_sig}, \code{signals}, \code{units},
\code{sig\_len}, \code{duration\_s}, \code{start\_utc}, \code{time\_source},
\code{time\_uncertainty\_s}, provenance, checksums) are for people and other programs: the
readers take the rate, the signals and the length from the headers.

\begin{table}[htbp]
\caption{\file{gaps.csv}: one row per break in the data, or per change that started a new
record.}
\label{tab:recordings-gaps}
\begin{tabularx}{\linewidth}{@{}P{4.4cm}L@{}}
\toprule
{\tabhead Column} & {\tabhead Meaning} \\
\midrule
\code{scope} & \code{all} (every channel), or the modality or channels it concerns \\
\code{reason} & Why: \code{signal\_loss}, \code{corrupt\_bytes}, \dots{} The reasons
  \code{channel\_absent}, \code{channel\_set\_change}, \code{label\_change} and
  \code{packet\_count\_change} mark where a record was split while the data stayed contiguous:
  they are not shown as breaks \\
\code{end\_before\_unix}, \code{start\_after\_unix} & The end of the data before, and the start
  of the data after, when known \\
\code{gap\_s} & The length of the break, when known (even if neither end has a stated time) \\
\code{gap\_uncertainty\_s}, \code{measured\_by}, \code{exact}, \code{notes} & How it was
  measured: \code{exact} is \code{True} when \code{gap\_s} is exact \\
\code{records\_before}, \code{records\_after} & The records on each side, separated by
  \code{;}, by name or path \\
\bottomrule
\end{tabularx}
\end{table}

\begin{table}[htbp]
\caption{\file{vitals.csv}: one row per value.}
\label{tab:recordings-vitals}
\begin{tabularx}{\linewidth}{@{}P{3.0cm}L@{}}
\toprule
{\tabhead Column} & {\tabhead Meaning} \\
\midrule
\code{time} & When the value was measured \\
\code{label} (else \code{name}) & What was measured: \code{HR}, \code{SpO2}, \code{NBP-S} \dots \\
\code{value} & The value (rows whose value is not a number are left out) \\
\code{unit} & Its unit \\
\code{validity} & Empty or \code{0}: valid; anything else: the row is left out \\
\code{source} & Where the value came from, when a recording has several (the monitor's
  numerics, an import) \\
\bottomrule
\end{tabularx}
\end{table}
A label measured by one source is named by its label. When a label has values from several
sources, the first source in the file keeps the plain label and the others are named with their
source: \code{HR}, \code{HR [import]}. Series of different sources are never merged.

\begin{table}[htbp]
\caption{\file{events.csv}: alarms and other events that have a start and, possibly, an end.}
\label{tab:recordings-events}
\begin{tabularx}{\linewidth}{@{}P{2.4cm}L@{}}
\toprule
{\tabhead Column} & {\tabhead Meaning} \\
\midrule
\code{start\_unix} & The start (required) \\
\code{end\_unix} & The end; empty: an instant \\
\code{level} & The severity: 1 advisory, 2 warning, 3 crisis (shown by color) \\
\code{text} & What is shown \\
\code{kind} & The kind of event (\code{alarm}, the system that raised it \dots); default
  \code{event} \\
\bottomrule
\end{tabularx}
\end{table}
Events are kept when their start lies within a day of the recording's data (a guard against zero
dates), sorted by start.

\paragraph{Markers.} A record may have a \file{.mrk} annotation file (WFDB's MIT format): marks
at samples, such as R waves or pacing spikes. An annotation's \code{chan} is the signal it belongs
to. A marker's name is its auxiliary text when it has one (\code{R\_WAVE}), else the WFDB
mnemonic of its code (\code{N}, \code{V}, \code{+} \dots). Markers are drawn on their signal and
can be listed by name.

\paragraph{How a collection is read.}
\begin{enumerate}
\item \textbf{Records} are taken in \file{index.csv} order, each signal a piece of its channel
  (modality, signal description, unit). An \file{index.csv} without a \code{record} column is
  refused (\code{index.csv has no record column}). A row whose \code{record} is empty, absolute
  or has a \code{..} component is left out (under its name, or \code{row N} when empty, counting
  the rows after the header from 1), with the reason \code{index.csv names a record outside its
  folder (an absolute name, or one with ..)}; a record that cannot be read is left out with the
  reason, as in a folder of records; the recording's notes say how many were left out, and a
  collection of which no record can be read is refused.
\item \textbf{Placement on the time axis.} A record with \code{start\_unix} starts there. A record
  without one follows the data before it, after the break \file{gaps.csv} states before it
  (\code{gap\_s} on a row of scope \code{all} whose \code{records\_after} names it), and records
  of the same \code{stretch} start together. Such records are drawn as \emph{time not stated}:
  their place is known relative to their neighbors, not on the clock.
\item \textbf{Breaks} (\file{gaps.csv}, except the record splits) are drawn between the data,
  with their reason and length; a break without stated times takes them from the records on each
  side.
\item \textbf{Vital signs, events and markers} as above.
\end{enumerate}

\paragraph{An example.} The smallest collection: one ECG record with a start time, and an alarm.
\begin{lstlisting}[style=text]
index.csv
  record,modality,start_unix
  ecg/rec_0001,ecg,1709287200.0
events.csv
  start_unix,end_unix,level,text,kind
  1709287201,1709287203,2,HR HIGH 151 > 150,alarm
ecg/RECORDS
  rec_0001
ecg/rec_0001.hea, ecg/rec_0001.dat
\end{lstlisting}
\file{viewers/vitalstream/testdata/generate.py} writes a complete synthetic one (three stretches,
a break of stated length, a record without a start time, vital signs from two sources, alarms and
markers) with the Python standard library only; \file{viewers/vitalstream/testdata/collection}
is the test collection the viewers and Medorama are checked on.

\subsection{Overlapping records and lanes}
\label{sec:recordings-overlap}

A signal recorded twice at the same time (two records of the same channel whose times overlap)
is never drawn on top of itself: each piece goes into the first \emph{lane} of its channel that
it does not overlap, and each further lane is a channel of its own, named with its lane:
\code{II}, \code{II (2)}, \code{II (3)} \dots{} Pieces are placed in the order of the records
(\file{index.csv}'s order, or their paths), and the signals of a record in header order. The
recording's notes say so when it happens. A channel is a modality (\code{wfdb} for plain
records), a signal description and a unit.

\subsection{Limits}
\label{sec:recordings-limits}

So that opening a large folder by mistake cannot overwhelm a viewer, a folder of plain records is
read only within the limits of \cref{tab:recordings-limits}. They are VitalStream settings
(\code{max\_records}, \code{max\_depth}, \code{max\_scan\_s} in the viewer's settings file;
\cref{chap:vitalstream}); Medorama has no such settings, and indexes a recording's folder within
the defaults.

\begin{table}[htbp]
\caption{The limits on a folder of plain records.}
\label{tab:recordings-limits}
\begin{tabularx}{\linewidth}{@{}LP{1.6cm}P{1.8cm}@{}}
\toprule
{\tabhead Limit} & {\tabhead Default} & {\tabhead Over it} \\
\midrule
Header files below the folder (segments and layout headers counted) & 5000 & Refused \\
Folder levels below the folder & 20 & Refused \\
Seconds to count the records before opening (the viewers' summary) & 10 & Refused \\
\bottomrule
\end{tabularx}
\end{table}
A refusal says which limit was reached and what to do instead: open a subfolder, or describe the
records in an \file{index.csv} (a collection, which lists its records itself and is not limited
by these). Every recording, a collection included, has at most 100\,000 pieces (signals
$\times$ records), and a folder's index key is computed from at most 200\,000 of its files
within \code{max\_depth} levels (\cref{sec:formats-reference-limits}).
Before opening a folder, the viewers show what it holds (records, segments, signals, time span,
size), read from the headers alone; indexing starts only when it is confirmed, shows its
progress, and can be cancelled.

\section{EDF, EDF+, BDF and BDF+ files}
\label{sec:recordings-edf}

The European Data Format (EDF) is a widely used open format for physiological recordings: sleep
studies (polysomnography), EEG, long-term ECG and many research datasets. Medorama, the three
VitalStream viewers and the ECG viewer read four variants directly, without an add-on
(\cref{tab:recordings-edf-variants}). The specifications are published at
\url{https://www.edfplus.info/specs/}.

\begin{table}[htbp]
\caption{The variants of EDF that Medorama and its viewers read.}
\label{tab:recordings-edf-variants}
\begin{tabularx}{\linewidth}{@{}P{1.6cm}P{1.6cm}L@{}}
\toprule
{\tabhead Variant} & {\tabhead Samples} & {\tabhead What it adds} \\
\midrule
EDF & 16-bit & The original format (Kemp et al., 1992). \\
EDF+ & 16-bit & Annotations, and recordings with interruptions (Kemp and Olivan, 2003). \\
BDF & 24-bit & The BioSemi variant of EDF, used by high-resolution EEG systems. \\
BDF+ & 24-bit & BDF with the EDF+ extensions. \\
\bottomrule
\end{tabularx}
\end{table}

An EDF file is one recording. A file is read as EDF when its name ends in \file{.edf} or
\file{.bdf}, in any letter case. The built-in reader takes precedence over any adapter that lists
these extensions. Medorama and the three viewers read the same file to the same index, byte for
byte. In VitalStream, an EDF file becomes the following:
\begin{itemize}
\item Each ordinary signal becomes a channel, with the signal's label and physical dimension
  (\code{ECG}, \code{mV}). Signals may have different sampling rates. A second signal with the
  same label becomes a channel named with its count, such as \code{ECG (2)}.
\item The header's start date and time give the time axis. The time is taken as UTC, since EDF
  does not state a time zone.
\item EDF+ and BDF+ annotations become events in the \ui{Annotations} lane, with their text and
  duration.
\item An interrupted EDF+ or BDF+ recording (EDF+D) gives one stretch of data per uninterrupted
  run, with a break between runs (``not recorded'').
\end{itemize}
The exact reading rules, which the three implementations follow, are specified in
\cref{sec:formats-reference-edf}; what the viewer shows is described in
\cref{sec:vitalstream-edf}.

\paragraph{EEG and sleep studies in Medorama.} Medorama has two recording categories for these
files, \textbf{EEG} (\code{eeg}) and \textbf{Sleep} (\code{sleep}), beside ICU and Holter. Each is
a column of the patient page's journey, and its recordings open in VitalStream. A table's
\code{category} column, the category chosen for the import, or the table's file name (with words
such as \code{eeg}, \code{polysomnography}, \code{psg} or \code{sleep study}) chooses the
category (\cref{chap:importing}). With the category \code{ecg}, an EDF file is not a recording
but an attachment of its entry, opened by the ECG viewer (\cref{chap:ecg-imaging}).

The demo made by \file{tools/make-demo-data.py} includes an overnight sleep study for the
patient \code{P900004}: a two-hour EDF+ file (\file{sleep/P900004\_psg\_20240305.edf}) with EEG,
EOG, ECG, airflow, thoracic effort and SpO2 signals, and its scoring (sleep stages and events) as
annotations. The study has its own encounter in the sleep laboratory (\code{E9100010}), with a
note, the \emph{Sleep study report}. Its table, \file{sleep\_index.csv}, links the file to the
patient and the encounter:
\begin{lstlisting}[style=text]
patient_id,encounter_id,time,study,category,file
P900004,E9100010,2024-03-05 22:30,Polysomnography,sleep,../sleep/P900004_psg_20240305.edf
\end{lstlisting}

\screenshot{patient-sleep.png}{The demo's EDF+ sleep study in the Sleep column of a patient's
page (dark theme).}{fig:recordings-sleep}

\paragraph{The patient's identity in the header.} An EDF header has two identification fields of
80 characters each. EDF+ structures the \emph{local patient identification} as
\code{code sex birthdate name}, with \code{X} for any part that is unknown, and the \emph{local
recording identification} as \code{Startdate dd-MMM-yyyy admincode technician equipment}. These
fields often hold protected health information, so they are handled as follows:
\begin{itemize}
\item The index never contains them, and neither does anything the viewers send to the browser.
  VitalStream and the ECG viewer do not show them, in Medorama or on their own.
\item Medorama's \code{identify} action reads them (\code{from\_edf} in
  \file{medorama/src/identify.rs}). It stores the code, sex, birth date, name and admin code in
  the catalog's identifier tables, next to the database's own values, so that an administrator
  can check that the file belongs to the patient it is linked to
  (\cref{sec:recordings-identifiers}). A plain EDF patient field, which has no structure, is kept
  whole as the code.
\item In Privacy Mode, EDF files and their annotations are de-identified before they leave the
  server (\cref{sec:protecting-edf}).
\end{itemize}

\paragraph{Limits.}
\begin{itemize}
\item \textbf{A folder of EDF files is not yet one session.} VitalStream's \ui{Open multiple
  records\dots} reads the WFDB records of a folder; open EDF files one at a time. The ECG viewer's
  \ui{Open multiple records\dots} does list the EDF files of a folder.
\item \textbf{Annotation-only files are refused.} An EDF+ file that holds only annotations (some
  sleep-scoring files do) has no signal to draw. Its events cannot yet be shown over another
  file's signals.
\end{itemize}

\section{Formats read by add-ons}
\label{sec:recordings-addons}

Medorama and the viewers read WFDB records and EDF files themselves. Every other recording
format is read by an \emph{adapter}: a separate program, installed beside them in the adapter
folder (the setting \code{adapters\_dir}), that turns a recording into the index the viewers draw
from, and reads its samples when a view zooms in. Formats that are licensed, or that only some
institutions may read, are distributed this way, as \emph{add-ons}, outside the repository. An
adapter writes only the index files, into the folder it is given, and never anything next to a
recording. It is called as a child process, and listens on no port.

\paragraph{BedMasterEx \file{.Stp}.} The STP add-on (\code{stp-adapter}) reads the \file{.Stp}
files of BedMasterEx: recordings of GE monitors (Unity and Carescape) and of Philips monitors
(PIICiX and Viridia). It is private: it is licensed separately and released, as a prebuilt
program, to institutions under agreement. It is not in Medorama's repository.

Until an add-on is installed, Medorama and the viewers name none of its formats: a file of that
format is treated like any other file they do not read. The viewers' file browsers do not list
it, and opening it is refused (\code{not a recording this program reads}). An import skips the row
that names it, with a warning, and keeps nothing of it. Once an administrator has installed the
add-on on the server (never through the page) and restarted Medorama, remove that earlier import
and import the table again. The files then come in as recordings, linked to their patients and
indexed by the workers. Installing, checking and removing add-ons, and their
protocol, are described in \cref{chap:addons}.

With an adapter, the recording is indexed like any other: the index holds its channels, pieces,
vital signs, alarms, breaks and markers, the format's label, and the adapter's name. Raw samples
are read from the recording through the adapter when a view zooms in close enough. Indexes made
through an adapter are shared like the others: the three viewers and Medorama read and write the
same index folders. A failing adapter makes the recording unreadable, never the program; the
indexing fails with the adapter's message.

\section{Linking recordings to encounters}
\label{sec:recordings-linking}

A recording is linked to one or more encounters after every import and every indexing, in this
order:
\begin{enumerate}
\item \textbf{By a database row} that names its file and its encounter (the \code{file} and
  \code{encounter\_id} columns). The recording belongs to that encounter, as if linked manually.
\item \textbf{By a bed stay's \code{recording\_file}.} A recording whose file name or path is a
  bed stay's \code{recording\_file} (ignoring case) is linked to that bed stay's encounter.
\item \textbf{By bed and time.} Bedside-monitor systems name their files
  \code{UNIT\_BED-<Unix time>} (\file{ICU7\_12-1704106800.Stp}): the unit is the part of the name
  before the first \code{\_}, the bed the rest before the last \code{-}, and the time the digits
  after it (9 to 11 digits). A recording so named is linked to every encounter that has a bed stay
  whose \code{care\_unit} and \code{bed} are its unit and bed (\code{care\_unit} + \code{\_} +
  \code{bed}, ignoring case) and whose time overlaps the recording's. Before indexing, the
  recording's time is the time in its name; after indexing, it is the real start and end. A bed
  stay without an end is open.
\end{enumerate}

\paragraph{Decisions made manually} win over all of these, and outlive removing and importing the
databases again. In an encounter's panel (\cref{chap:patients}), admins have, for each recording:
\begin{itemize}
\item \ui{Unlink} unlinks the recording, and it is \emph{never} linked to that encounter
  automatically again.
\item \ui{Forget manual link}, for a recording linked manually, forgets the decision, and
  automatic linking resumes.
\end{itemize}
A recording with a manual link gets exactly its manual links. For administrators with the PHI
permission, \ui{Review} lists the \ui{Recordings whose files disagree}: automatic links where the
recording's own file names a patient other than the one in the databases. \ui{Keep the link}
makes such a link manual; \ui{Unlink} never links the recording to that encounter again. Every decision is in the audit log.
Recordings that match nothing are listed in \ui{Import} $\rightarrow$ \ui{Recordings} with
\ui{Not linked to an encounter} (\cmd{medorama recordings -{}-unlinked}).

\section{Indexing in the background}
\label{sec:recordings-indexing}

A recording is drawn from its \emph{index}: a folder in \code{index\_dir} that holds its
channels, the envelopes of every piece (the minimum and maximum of every block of 32 samples, and
of every four blocks above that), its vital signs, alarms, breaks and markers. The folder's name is
the recording's \emph{index key}: a fingerprint of the index version, the recording's path and
its files' names, sizes and modification times. A changed recording has a new key, and is indexed
again.

\begin{itemize}
\item \textbf{Automatically.} With \code{auto\_index} on (the default), the server queues the
  indexing of every recording that is new or changed, and the workers (\code{workers}, default 2)
  index them one by one while it runs. The \ui{Jobs} tab shows the queue (\cref{chap:actions}).
  Users who may run actions can also \ui{Index now} a recording that is not indexed.
\item \textbf{Shared.} A recording whose current index is already in \code{index\_dir} (for
  instance, the viewers' \code{cache\_dir}, pointed at the same folder) is marked indexed at once,
  without indexing it again.
\item \textbf{Times.} Once a recording is indexed, its start and end are its real ones, and its
  format is the index's (\code{WFDB records}, \code{EDF+}, or the add-on's label). Links by bed
  and time are computed again.
\item \textbf{Identifiers.} The built-in action \code{identify} reads the identifiers that the
  recording states, once per recording and again when its file changes. Medorama reads an EDF
  header's fields itself, and asks an adapter for the formats it reads; WFDB records hold none
  (\cref{sec:recordings-identifiers}).
\end{itemize}

\paragraph{Checking the files.} \ui{Re-check the files} (\ui{Import} $\rightarrow$
\ui{Recordings}), \cmd{medorama check}, and the server every \code{check\_interval\_s} seconds
(default 900; 0 means never) check every recording: whether it is still there, whether it
changed, and whether its index folder is still there. A missing recording is kept, marked
missing, until an administrator removes it, and is counted as back when it returns. A changed
recording is indexed again, and a lost index is made again. The following command runs the check
and prints a summary:
\begin{lstlisting}[style=shell]
medorama check
# 1204 checked: 2 missing, 0 back, 3 changed, 0 index folders gone
\end{lstlisting}

\section{The Recordings tab}
\label{sec:recordings-tab}

\ui{Import} $\rightarrow$ \ui{Recordings}, for administrators, lists every recording, a page at
a time. For each recording it shows the status, the name (with any error, and, for users with the
PHI permission, what its file says of the patient), the kind and format, the start, the length,
the encounters, the sources and the path, with the buttons \ui{View} (in VitalStream),
\ui{Copy path}, \ui{Index now} and \ui{Remove}. The status is \code{new} (never indexed),
\code{queued}, \code{indexing}, \code{ready}, \code{stale} (changed since it was indexed),
\code{failed} or \code{missing} (not found at the last check). The search line takes words or
\code{field:value} terms (\code{file:}, \code{path:}, \code{unit:}, \code{source:},
\code{status:} \dots). \ui{Not linked to an encounter} keeps only the unlinked recordings, and
\ui{Remove matching\dots} removes every recording the search finds, after saying how many.

\paragraph{Removing recordings} takes them out of the catalog; the files are never touched. A
recording's index and results are deleted too only when asked (the page asks; on the command
line, \opt{-{}-purge-results}), and only inside \code{results\_dir} and \code{index\_dir}.
Encounters stay. The following commands list the unlinked recordings, show what a removal would
take out, and remove it:
\begin{lstlisting}[style=shell]
medorama recordings --unlinked                       # recordings no encounter holds
medorama remove --search "unit:ICU7" --dry-run       # what would go
medorama remove --search "unit:ICU7" --purge-results
\end{lstlisting}

\section{Identifiers in recordings}
\label{sec:recordings-identifiers}

\begin{phibox}
Bedside monitors and EDF headers store patient identifiers in the recordings themselves. The
action \code{identify} reads them: from an EDF header's patient and recording fields itself, and
through the recording's adapter for other formats. A file that names several patients is
flagged. The identifiers are kept in the catalog's database only, and shown only to users with
the PHI permission outside Privacy Mode, so that a reviewer can check the databases against them.
They appear in the Encounters table's \ui{In the files}, the encounter panel, \ui{Import}
$\rightarrow$ \ui{Recordings}, and the searches \code{file-name:}, \code{file-id:} and
\code{file-dob:}. They never reach job summaries, the viewers, actions' job descriptions,
results folders or logs.

The index holds no name, identifier or date of birth. A recording's path and file name may still
identify a patient, however (bedside-monitor systems name files by unit, bed and time), and every
value of a recording is patient data. Keep recordings, and the index folders made from them, on
protected storage. In Privacy Mode, VitalStream opens recordings by token, never by path, and
every time it shows is moved (\cref{sec:protecting-vitalstream}).
\end{phibox}
